% Balance the references in this draft; adjust or remove after repagination.
\IEEEtriggeratref{6}
\begin{thebibliography}{00}

\bibitem{qvhighlights}
J. Lei, T. L. Berg, and M. Bansal, ``QVHighlights: Detecting moments and highlights in videos via natural language queries,'' in \emph{Adv. Neural Inf. Process. Syst.}, vol. 34, 2021, pp. 11846--11858.

\bibitem{nextqa}
J. Xiao, X. Shang, A. Yao, and T.-S. Chua, ``NExT-QA: Next phase of question-answering to explaining temporal actions,'' in \emph{Proc. IEEE/CVF Conf. Comput. Vis. Pattern Recognit.}, 2021, pp. 9777--9786.

\bibitem{egoschema}
K. Mangalam, R. Akshulakov, and J. Malik, ``EgoSchema: A diagnostic benchmark for very long-form video language understanding,'' in \emph{Adv. Neural Inf. Process. Syst.}, vol. 36, 2023.

\bibitem{lewis}
P. Lewis \emph{et al.}, ``Retrieval-augmented generation for knowledge-intensive NLP tasks,'' in \emph{Adv. Neural Inf. Process. Syst.}, vol. 33, 2020, pp. 9459--9474.

\bibitem{videorag}
Y. Luo \emph{et al.}, ``Video-RAG: Visually-aligned retrieval-augmented long video comprehension,'' in \emph{Adv. Neural Inf. Process. Syst.}, vol. 38, 2025.

\bibitem{whisper}
A. Radford \emph{et al.}, ``Robust speech recognition via large-scale weak supervision,'' in \emph{Proc. Int. Conf. Mach. Learn.}, vol. 202, 2023, pp. 28492--28518.

\bibitem{pyannote}
H. Bredin \emph{et al.}, ``pyannote.audio: Neural building blocks for speaker diarization,'' in \emph{Proc. IEEE Int. Conf. Acoust. Speech Signal Process.}, 2020, pp. 7124--7128.

\bibitem{rrf}
G. V. Cormack, C. L. A. Clarke, and S. B\"uttcher, ``Reciprocal rank fusion outperforms Condorcet and individual rank learning methods,'' in \emph{Proc. Int. ACM SIGIR Conf. Res. Develop. Inf. Retrieval}, 2009, pp. 758--759, doi: 10.1145/1571941.1572114.

\bibitem{legacy}
A. Anish and V. Verma, ``FalconVQA,'' project development report, \emph{Product Performance} and \emph{Dataset Creation} sections, repository file \texttt{latex/phases.tex}, rev. \texttt{76fc7ff}.

\end{thebibliography}
